\documentclass[preprint,12pt,authoryear]{elsarticle}
\usepackage{amsmath,amssymb,amsthm,booktabs,graphicx,tabularx,float}
\usepackage[T1]{fontenc}
\usepackage{lmodern}
\usepackage[utf8]{inputenc}
\usepackage[hidelinks]{hyperref}
\usepackage{placeins}
\pdfstringdefDisableCommands{%
  \def\corref#1{}%
}
\usepackage{longtable}
\newtheorem{proposition}{Proposition}
\newtheorem{theorem}{Theorem}
\newtheorem{lemma}{Lemma}
\newcounter{algorithm}
\newenvironment{boxedalgorithm}[1]{%
  \refstepcounter{algorithm}%
  \par\medskip
  \noindent
  \begin{minipage}{\textwidth}
  \hrule\vspace{1pt}
  \textbf{Algorithm \thealgorithm. #1}\par\vspace{1pt}
  \hrule\vspace{2pt}\fontsize{9.0}{9.6}\selectfont
}{%
  \hrule
  \end{minipage}
  \par\medskip
}
\journal{Computational Statistics \& Data Analysis}

\begin{document}
\begin{frontmatter}
\title{Centered Natural-Parameter Group Selection for Posterior-Score Contrast Support in von Mises--Fisher Mixtures}
\author[aff1]{Hyunsoo Shin}
\author[aff1]{Byungtae Seo\corref{cor1}}
\cortext[cor1]{Corresponding author.}
\address[aff1]{Department of Statistics, Sungkyunkwan University, Seoul 03063, Republic of Korea}
\ead{seobt@skku.edu}
\begin{abstract}
\end{abstract}
\begin{keyword}Directional data; mixture model; von Mises--Fisher distribution; heterogeneity selection; group lasso; posterior-score contrast support\end{keyword}
\end{frontmatter}

\section{Introduction}

Directional observations include normalized documents, compositional
profiles, and hyperspherical embeddings. Finite von Mises--Fisher (vMF)
mixtures provide a model-based alternative to spherical $k$-means
\citep{banerjee2005vmf}, but high-dimensional fits can be difficult to estimate
and interpret.

Sparse vMF mixtures have primarily penalized component directions to estimate
sparse prototypes \citep{rossi2022sparsevmf}. Prototype presence is not the
same as contribution to posterior component comparisons: a coordinate can be
active in every prototype yet cancel from score contrasts, whereas unequal
concentrations can amplify modest directional differences.

Related mixture methods use entrywise shrinkage, grouped means, pairwise
fusion, or discriminative subspaces
\citep{pan2007penalized,xie2008grouped,guo2010pairwise,
bouveyron2012fisher,bouveyron2014sparsefisher}; sum-to-zero contrasts separate
common levels from heterogeneity \citep{bondell2009anova}. Directional-mixture
alternatives target saliency, blocks, or hierarchical regularization
\citep{amayri2013online,salah2016coclustering,gopal2014vmf,
taghia2014bayesian,salah2019directional}. Here component-specific concentration
makes natural-parameter, rather than directional, contrasts relevant to
posterior scores.

Separating support construction from likelihood refitting parallels
Lasso--MLE model collections for Gaussian mixtures
\citep{meynet2012thesis,meynet2012sparse}. Our construction instead aligns
both stages with centered vMF natural-parameter support, which equals the
support of pairwise posterior-score linear coefficients. The common baseline
cancels from those coefficients but remains in the fitted density.

This paper targets posterior-score contrast support through natural parameters
formed by multiplying each component direction by its concentration. For each
observed coordinate, the component-specific natural parameters are decomposed
into a common baseline and a centered contrast. The proposed centered
natural-parameter group lasso (E-CGL) penalty
groups the centered contrast across components, so its unit of regularization
is the same coordinate-level unit used to define the target support. The common
baseline remains unpenalized. An adaptive centered natural-parameter group
lasso (E-ACGL) is retained as a secondary extension; E-CGL remains the primary
estimator.

Candidate supports are generated by a guarded path with proximal-gradient
updates inside a generalized expectation-maximization (GEM) algorithm. They
are compared after support-constrained
refitting that preserves the target.
The primary selector is a practical Bayesian information criterion (BIC)
computed from the refitted likelihood; alternative degrees of freedom and
extended BIC (EBIC) rules are treated as sensitivity
analyses. Support-recovery claims condition on a fixed number of components.
When that number is unknown, component-number selection and support selection
are examined in separate stages rather than through a single joint criterion.

The empirical study has three parts. Oracle-error calibrated simulations
compare posterior-score contrast support recovery across sample size, ambient
dimension, and concentration heterogeneity. Shared-natural-parameter and
dense-support designs separate the proposed estimand from prototype sparsity
and identify settings in which assuming sparse posterior-score contrast
support is inappropriate. A full-data Classic3
analysis examines clustering, coordinate reduction, and token-level contrasts
on a vocabulary-aligned directional representation. Prespecified Classic3
splits assess support stability.

The contributions are therefore as follows: (i) a posterior-score contrast support
estimand for vMF mixtures, its centered representation, and its decomposition
into directional and concentration effects; (ii) E-CGL, together with guarded
estimation and target-preserving refitting, with E-ACGL as an adaptive
extension using fixed pilot weights;
(iii) existence on the bounded computational parameter space, convergence of
the fixed-responsibility proximal algorithm, and conditional stationarity of
accumulation points under stated regularity conditions; and
(iv) an empirical evaluation that documents sparse, prototype, and
dense-support regimes.
E-CGL is not intended to replace sparse-prototype methods when prototype
sparsity is the scientific target.

The remainder of the paper is organized as follows. Section 2 presents the vMF
model, posterior-score contrast support, and centered group regularization. Section 3
develops the guarded proximal GEM estimator, pathwise support selection and
refitting, and the
structural and algorithmic properties used in the analysis. Sections 4 and 5
report the numerical studies and real-data analyses, respectively; the latter
use prespecified preprocessing rules and split indices.
Section 6 summarizes the main findings, limitations, and conclusions.

\section{Model and centered regularization}

\subsection{vMF mixture in natural parameters}

Let $x_i\in\mathbb S^{d-1}$, $i=1,\ldots,n$, denote unit-normalized
directional observations. A $K$-component vMF mixture has density

\[
f(x_i;\Theta)
=\sum_{k=1}^K\pi_k C_d(\kappa_k)
\exp\{\kappa_k\mu_k^\top x_i\},
\]

where $\pi_k>0$, $\sum_k\pi_k=1$, $\lVert\mu_k\rVert_2=1$, and
$\kappa_k\geq0$. Its normalizing constant is

\[
C_d(\kappa)
=\frac{\kappa^{d/2-1}}
{(2\pi)^{d/2}I_{d/2-1}(\kappa)}.
\]

Here $I_\nu$ is the modified Bessel function of the first kind of order $\nu$.

At $\kappa=0$, this expression is interpreted by continuous extension as
$C_d(0)=\Gamma(d/2)/(2\pi^{d/2})$.

Write
\[
\begin{aligned}
\Theta&=(\pi_1,\ldots,\pi_K,\eta_1,\ldots,\eta_K),\\
E&=(\eta_1,\ldots,\eta_K)^\top,
&M&=(\mu_1,\ldots,\mu_K)^\top,
\end{aligned}
\]
where $E,M\in\mathbb R^{K\times d}$ and
$\eta_k=\kappa_k\mu_k\in\mathbb R^d$. Then

\[
f(x_i;\Theta)
=\sum_{k=1}^K\pi_k C_d(\lVert\eta_k\rVert_2)
\exp(\eta_k^\top x_i).
\]

The component-level mapping back to $(\mu_k,\kappa_k)$ and its
zero-concentration exception are recorded in Section~\ref{sec:structural}. The
observed log-likelihood is

\[
\ell(\Theta)
=\sum_{i=1}^n\log\left\{
\sum_{k=1}^K\pi_k C_d(\lVert\eta_k\rVert_2)
\exp(\eta_k^\top x_i)\right\},
\]

where the component densities are defined with respect to surface-area
measure on $\mathbb S^{d-1}$.

\subsection{Posterior-score contrast support}

Let $Z\in\{1,\ldots,K\}$ denote the latent component label for a generic
observation $x$. The component score

\[
s_k(x)=\log\pi_k+\log C_d(\lVert\eta_k\rVert_2)+\eta_k^\top x
\]

implies the pairwise posterior log-odds representation

\[
\log\frac{\Pr(Z=k\mid x)}{\Pr(Z=\ell\mid x)}
=b_{k\ell}+(\eta_k-\eta_\ell)^\top x,
\]

\[
b_{k\ell}
=\log\frac{\pi_k}{\pi_\ell}
+\log\frac{C_d(\lVert\eta_k\rVert_2)}
{C_d(\lVert\eta_\ell\rVert_2)}.
\]

Thus $\eta_k-\eta_\ell$, rather than $\mu_k-\mu_\ell$, is the coefficient
of the feature-dependent linear term. Let $\eta_{\cdot j}$ denote column $j$
of $E$, and define

\[
H_K=I_K-K^{-1}\mathbf1_K\mathbf1_K^\top,
\qquad
c_{\cdot j}=H_K\eta_{\cdot j}
=\eta_{\cdot j}-\bar\eta_j\mathbf1_K,
\]

where $\bar\eta_j=K^{-1}\sum_k\eta_{kj}$. We define

\[
S_{\eta}
=\{j:\lVert c_{\cdot j}\rVert_2>0\}.
\]

For population natural parameters $E^\star$, write
$S_\eta^\star=\{j:\lVert H_KE^\star_{\cdot j}\rVert_2>0\}$.

We use \emph{posterior-score contrast support} for this feature-dependent
linear support. It does not include the intercepts
$b_{k\ell}$. A coordinate with $c_{\cdot j}=0$ cancels from every pairwise
linear term, although its common baseline $\bar\eta_j$ remains in the fitted
component densities and may affect $b_{k\ell}$ through the concentration
norms. Proposition~\ref{prop:centered-support} shows that this support is the
union of the feature-dependent linear supports in all pairwise posterior
log-odds and gives the corresponding pairwise-dispersion identity.

For directional-support comparisons, assume $\kappa_k>0$ for every component,
so that $\mu_k=\eta_k/\lVert\eta_k\rVert_2$ is identified. Then let
\[
\bar\mu_j=K^{-1}\sum_{k=1}^{K}\mu_{kj},
\qquad
c_{\cdot j}^{(\mu)}
=H_KM_{\cdot j}
=M_{\cdot j}-\bar\mu_j\mathbf 1_K,
\]
and define
\[
S_P=\{j:\lVert M_{\cdot j}\rVert_2>0\},
\qquad
S_\mu=\{j:\lVert c_{\cdot j}^{(\mu)}\rVert_2>0\}.
\]
The set $S_P$ records coordinates present in at least one directional
prototype, whereas $S_\mu$ records between-component variation in mean
directions. The proposed $S_\eta$ records coordinates entering
posterior-score contrasts. The three targets can differ when component
concentrations are heterogeneous. If any $\kappa_k=0$, the corresponding
direction is unidentified and $S_P$ and $S_\mu$ require an explicit directional
convention; $S_\eta$ remains well-defined through the natural parameters.

Accordingly, $c_{\cdot j}\neq0$ defines a posterior-score contrast coordinate;
$c_{\cdot j}=0$ with $\bar\eta_j\neq0$ defines a common coordinate; and
$c_{\cdot j}=0$ with $\bar\eta_j=0$ defines a globally null coordinate.
These categories distinguish the proposed target from prototype support,
which records nonzero component coordinates without requiring a difference
between components.

The support is invariant to permutations of component labels because such a
permutation only reorders the entries of $c_{\cdot j}$. It is not invariant to
an arbitrary orthogonal rotation of the feature axes. As with other
coordinate-selection procedures, its interpretation is therefore relative to
the chosen coordinate system.

\subsection{Centered natural-parameter group lasso (E-CGL)}

The proposed E-CGL method regularizes the vMF natural parameters
$\eta_k=\kappa_k\mu_k$, collected row-wise in
$E\in\mathbb R^{K\times d}$. It penalizes the coordinatewise contrasts
$H_KE_{\cdot j}$ while leaving the common natural-parameter baseline
unpenalized. For fixed $K$, the primary penalty functional is

\[
\mathcal P_{\mathrm{CGL}}(E)
=\sum_{j=1}^d
\lVert H_KE_{\cdot j}\rVert_2
=\sum_{j=1}^d
\lVert c_{\cdot j}\rVert_2.
\]

Each group is one observed coordinate across the $K$ component contrasts, so
the unit of regularization matches the definition of $S_{\eta}$.
This coordinate-group construction follows the grouped-variable
regularization principle of \citet{yuan2006group}.
By contrast,

\[
\lVert\eta_{\cdot j}\rVert_2^2
=K\bar\eta_j^2+\lVert c_{\cdot j}\rVert_2^2,
\]

so a raw-$\eta$ group penalty combines common and contrast effects. A
centered direction-space penalty targets $S_\mu$, whereas the Rossi-style
entrywise penalty targets $S_P$. When the $\kappa_k$ differ, neither target is
generally equal to the coefficient support of $\eta_k-\eta_\ell$ in posterior
log-odds. These are distinct estimands, not competing parameterizations of the
same support.

An entrywise centered penalty selects component-coordinate entries separately:
\(
\lambda_\eta\sum_{k,j}|c_{kj}|.
\)
E-CGL instead selects the complete coordinate contrast as one group. The
resulting proximal map and label invariance are stated in
Propositions~\ref{prop:label-invariance} and~\ref{prop:centered-prox}.

\subsection{Adaptive weighted extension (E-ACGL)}

E-ACGL applies adaptive penalization at the coordinate-group level using
weights computed once from a dense pilot, following adaptive lasso and
adaptive group-lasso principles \citep{zou2006adaptive,wang2008adaptivegroup}.
Let
$a\in\mathcal A=\{\mathrm{CGL},\mathrm{ACGL}\}$ denote the method and let
$E^{\mathrm D}$ be the retained dense pilot. Define

\begin{equation}
\label{eq:eacgl-penalty}
\begin{aligned}
w_j^{(a)}
&=
\begin{cases}
1,
&a=\mathrm{CGL},\\[4pt]
\displaystyle
\frac{
\left(\lVert H_KE_{\cdot j}^{\mathrm D}\rVert_2+\epsilon\right)^{-\gamma}
}{
\operatorname{median}_{1\leq h\leq d}
\left(\lVert H_KE_{\cdot h}^{\mathrm D}\rVert_2+\epsilon\right)^{-\gamma}
},
&a=\mathrm{ACGL},
\end{cases}\\
\mathcal P_a(E)
&=\sum_{j=1}^d
w_j^{(a)}\lVert H_KE_{\cdot j}\rVert_2.
\end{aligned}
\end{equation}

The two methods share the dense pilot. E-CGL sets $w_j^{(a)}=1$ for every
$j$, whereas E-ACGL uses fixed pilot weights with $\gamma=1$ and
$\epsilon=10^{-6}$. Both procedures target $S_\eta$. For finite positive
weights, the results below extend by
replacing $\lambda_\eta$ with $\lambda_\eta w_j^{(a)}$ in group $j$. No
adaptive oracle claim is made.

\section{Estimation by guarded proximal generalized EM}
\label{sec:estimation}

\subsection{Penalized objective}

The finite vMF mixture likelihood can be unbounded as a component
concentration diverges \citep{ng2023penalizedvmf}. We therefore define the
fitted model on
\[
\mathcal T_{\kappa_{\max}}
=\left\{(\pi,E):\pi\in[0,1]^K,\ \mathbf1_K^\top\pi=1,\
\lVert\eta_k\rVert_2\leq\kappa_{\max}\ \text{for every }k\right\},
\]
with $\kappa_{\max}=10^6$ in all reported analyses. The numerical mixing-weight
floor is an implementation safeguard; the existence argument below uses the
closed simplex, and the local convergence statement assumes the floor and the
concentration bound are inactive.

For fixed $K$, $\lambda_\eta$, and $a\in\mathcal A$, the path-generating
criterion is

\begin{equation}
\label{eq:penalized-objective}
\mathcal L_{\lambda_\eta}^{(a)}(\Theta)
=\ell(\Theta)-\lambda_\eta\mathcal P_a(E).
\end{equation}

The penalty and the vMF normalizing constant are coupled through $\eta_k$, so
the natural-parameter update is computed by a guarded
proximal-gradient generalized M-step.

\begin{lemma}[Existence on the finite parameter space]
\label{lem:finite-existence}
For fixed $K,d,n$, finite $\kappa_{\max}$, nonnegative $\lambda_\eta$, and
finite positive fixed weights, the criterion in
\eqref{eq:penalized-objective} attains a maximum on
$\mathcal T_{\kappa_{\max}}$.
\end{lemma}

The lemma establishes existence for the bounded computational problem. It
does not imply uniqueness, mixture identifiability, or attainment of a global
maximum by the numerical algorithm.

\subsection{Fixed-penalty guarded proximal GEM}
\label{sec:fixed-gem}

We use a generalized expectation-maximization (GEM) procedure at each fixed
value of $\lambda_\eta$. The E-step computes the posterior responsibilities
exactly. In the generalized M-step, the mixing proportions are updated by
\eqref{eq:mixing-update}, while the natural-parameter block receives finitely
many safeguarded proximal-gradient updates. This block is required to increase
the penalized auxiliary criterion within numerical tolerance rather than to
attain its exact conditional maximum. The procedure is therefore a GEM
algorithm rather than an exact EM algorithm.

\subsubsection{E-step and fixed-responsibility subproblem}
\label{sec:fixed-responsibility}

At outer iteration $t$, compute $\tau_{ik}^{(t)}$ and

\begin{equation}
\label{eq:responsibilities}
\tau_{ik}^{(t)}
=\frac{\pi_k^{(t)}C_d(\lVert\eta_k^{(t)}\rVert_2)
\exp\{(\eta_k^{(t)})^\top x_i\}}
{\sum_{h=1}^K\pi_h^{(t)}C_d(\lVert\eta_h^{(t)}\rVert_2)
\exp\{(\eta_h^{(t)})^\top x_i\}},
\end{equation}

and define

\[
N_k^{(t)}=\sum_{i=1}^n\tau_{ik}^{(t)},
\qquad
r_k^{(t)}=\sum_{i=1}^n\tau_{ik}^{(t)}x_i.
\]

Conditional on these responsibilities and the fixed weights for method $a$,
the cap-inactive natural-parameter block is represented by

\begin{equation}
\label{eq:fixed-responsibility-objective}
\begin{aligned}
\min_E\quad F_t^{(a)}(E)
&=f_t(E)+g_{\lambda_\eta}^{(a)}(E),\\
f_t(E)
&=-\sum_{k=1}^K
\left\{
(r_k^{(t)})^\top\eta_k
+N_k^{(t)}\log C_d(\lVert\eta_k\rVert_2)
\right\},\\
g_{\lambda_\eta}^{(a)}(E)
&=\lambda_\eta\sum_{j=1}^d
w_j^{(a)}\lVert H_K\eta_{\cdot j}\rVert_2.
\end{aligned}
\end{equation}

The numerical algorithm evaluates this objective through cap-feasible trial
iterates. It rejects rather than projects a proposal that violates the
concentration cap; the cap-active constrained subproblem is not claimed to be
solved by the closed-form proximal update.

The smooth term is convex because
$-\log C_d(\lVert\eta_k\rVert_2)$ is the vMF log-partition function. Its
gradient is

\begin{equation}
\label{eq:fixed-responsibility-gradient}
\nabla_{\eta_k}f_t(E)
=N_k^{(t)}A_d(\lVert\eta_k\rVert_2)
\frac{\eta_k}{\lVert\eta_k\rVert_2}-r_k^{(t)},
\end{equation}

with the continuous value $-r_k^{(t)}$ at $\eta_k=0$, where
\[
A_d(\kappa)=\frac{I_{d/2}(\kappa)}{I_{d/2-1}(\kappa)}
\]
is the vMF mean-resultant-length function, with $A_d(0)=0$ by continuous
extension.

\subsubsection{Generalized M-step and centered proximal update}
\label{sec:centered-pg-update}

Let $J_K=K^{-1}\mathbf1_K\mathbf1_K^\top$ and $H_K=I_K-J_K$.
For a step size $s>0$, set

\begin{equation}
\label{eq:proximal-gradient-trial}
V=E^{(u)}-s\nabla f_t(E^{(u)}).
\end{equation}

The update separates by observed coordinate and has the closed form

\begin{equation}
\label{eq:centered-prox-update}
E_{\cdot j}^{(u+1)}
=J_KV_{\cdot j}
+\left(
1-\frac{s\lambda_\eta w_j^{(a)}}
{\lVert H_KV_{\cdot j}\rVert_2}
\right)_+H_KV_{\cdot j}.
\end{equation}

Here $(z)_+=\max(z,0)$, and the second term is defined as zero when
$H_KV_{\cdot j}=0$.

The common baseline is therefore retained, whereas the centered contrast is
group-thresholded; this is a proximal group soft-thresholding operation
\citep{parikh2014proximal}. Mixing proportions are updated by

\begin{equation}
\label{eq:mixing-update}
\pi_k^{\mathrm{cand}}
=\frac{\max(N_k^{(t)}/n,10^{-12})}
{\sum_h\max(N_h^{(t)}/n,10^{-12})}.
\end{equation}

The initial step size is $s=1/\max_kN_k^{(t)}$. Define the majorization
residual

\begin{equation}
\label{eq:majorization-residual}
\begin{aligned}
R_{\mathrm{maj}}(E^+,E;s)
={}&f_t(E^+)-f_t(E)
-\langle\nabla f_t(E),E^+-E\rangle\\
&-\frac{\lVert E^+-E\rVert_F^2}{2s}.
\end{aligned}
\end{equation}

The smooth majorization condition is equivalent to
$R_{\mathrm{maj}}(E^+,E;s)\leq0$. Numerically, the update is accepted only when
$R_{\mathrm{maj}}(E^+,E;s)\leq\varepsilon_{\mathrm{maj}}$ and
$\max_k\lVert\eta_k^+\rVert_2\leq\kappa_{\max}$. Otherwise, the proposal is
not projected in this proximal loop; the step is halved and the proposal is
recomputed. If
no acceptable proposal is found in 40 trials, the current path point is marked
invalid. The path driver skips that tuning value, retries the next value from
the last accepted fit, and terminates the remaining path only after three
consecutive failures.

\subsubsection{Safeguards and stopping rules}

Let $Q(\Theta\mid\tau)$ be the expected complete-data log-likelihood, using
the standard convention $0\log 0=0$ for zero-weight terms. For an old iterate
$\Theta$ and candidate $\Theta^+$, define

\begin{equation}
\label{eq:acceptance-increments}
\begin{aligned}
Q_{\lambda_\eta}^{(a)}(\Theta\mid\tau)
&=Q(\Theta\mid\tau)-\lambda_\eta\mathcal P_a(E),\\
\Delta Q_{\lambda_\eta}^{(a)}
&=Q_{\lambda_\eta}^{(a)}(\Theta^+\mid\tau)
-Q_{\lambda_\eta}^{(a)}(\Theta\mid\tau),\\
\Delta\mathcal L_{\lambda_\eta}^{(a)}
&=\mathcal L_{\lambda_\eta}^{(a)}(\Theta^+)
-\mathcal L_{\lambda_\eta}^{(a)}(\Theta).
\end{aligned}
\end{equation}

The guarded generalized M-step accepts the candidate only if

\begin{equation}
\label{eq:outer-acceptance}
\Delta Q_{\lambda_\eta}^{(a)}\geq-\varepsilon_{\mathrm{acc}},
\qquad
\Delta\mathcal L_{\lambda_\eta}^{(a)}\geq-\varepsilon_{\mathrm{acc}}.
\end{equation}

Otherwise, the path fit returns only the previous accepted estimate,
accompanied by a failure diagnostic; its support is not recorded, and the
path driver proceeds to the next tuning value from the last accepted fit.
Three consecutive failed tuning values terminate the remaining path.

Components with an effective count below $10^{-8}$ are retained to preserve
the fixed value of $K$, and this event is recorded. They are not pruned or
reinitialized. These safeguards support statements about accepted numerical
updates; they do not establish convergence of the complete parameter sequence,
stationarity of every fitted path point, or global optimality.

The recorded analyses use the absolute support threshold
$\varepsilon_{\mathrm{supp}}=10^{-8}$. This threshold is a numerical convention
for reading a fitted path support, not an additional penalty parameter.

The recorded simulation and real-data controls use at most 200 and 160
proximal-gradient iterations per generalized M-step, respectively. The inner
and outer stopping rules are

\begin{equation}
\label{eq:pg-stopping}
\frac{\lVert E^{(u+1)}-E^{(u)}\rVert_F}
{\max\{1,\lVert E^{(u)}\rVert_F\}}
<\varepsilon_{\mathrm{PG}},
\end{equation}

\begin{equation}
\label{eq:gem-stopping}
\frac{|\mathcal L_{\lambda_\eta}^{(a)}(\Theta^{(t+1)})
-\mathcal L_{\lambda_\eta}^{(a)}(\Theta^{(t)})|}
{\max\{1,|\mathcal L_{\lambda_\eta}^{(a)}(\Theta^{(t)})|\}}
<\varepsilon_{\mathrm{conv}}.
\end{equation}

The numerical tolerances are

\[
(\varepsilon_{\mathrm{PG}},\varepsilon_{\mathrm{maj}},
\varepsilon_{\mathrm{acc}},\varepsilon_{\mathrm{conv}},
\varepsilon_{\mathrm{supp}})
=(10^{-8},10^{-10},10^{-8},10^{-7},10^{-8}).
\]

The implementation permits at most 40 backtracking attempts and initially
allows 100 outer GEM iterations at each path point. If the provisional
refit-BIC winner is initialized by an unconverged penalized fit, the complete
path is rebuilt with 440 additional iterations per point and only converged
penalized sources are admitted to the repeated refit-BIC comparison.

\subsection{Pathwise selection and support-constrained refit}

\subsubsection{Method-specific zero-augmented KKT-motivated path}

Among multistart fits that are finite and free of numerical failure, the fit
with the largest observed log-likelihood is the provisional dense pilot. If it
reaches its iteration limit, its terminal
iterate is continued once for the same iteration budget. If continuation fails
or remains unconverged, the procedure falls back to the largest-likelihood
converged original start when available; otherwise the better finite terminal
fit is retained with a fallback flag. Denote the result by
$\widehat\Theta^{\mathrm D}$. After the weights in
\eqref{eq:eacgl-penalty} are computed and held fixed, the implementation first
obtains the guarded zero-penalty fit

\[
\widehat\Theta_0^{(a)}
=\operatorname{Fit}_{\lambda_\eta=0}^{(a)}
\left(\widehat\Theta^{\mathrm D};w^{(a)}\right).
\]

Responsibilities $\tau^0$ under $\widehat\Theta_0^{(a)}$ define
$(N_k^0,r_k^0)$ for path construction. Let
$r_\bullet=\sum_k r_k^0$, $N_\bullet=\sum_kN_k^0=n$, and
$\bar\rho_\bullet=\lVert r_\bullet\rVert_2/N_\bullet$. The Banerjee
approximation \citep[Eq.~(4.4)]{banerjee2005vmf} is

\begin{equation}
\label{eq:banerjee-pooled}
\widehat\kappa_B(\rho)
=\frac{\rho(d-\rho^2)}{1-\rho^2}
\end{equation}

Let $\widetilde\rho_\bullet$ denote the numerically clipped version of
$\bar\rho_\bullet$. The pooled natural parameter used for path construction is

\[
\eta^{\mathrm{common}}
=
\begin{cases}
\min\{\kappa_{\max},\widehat\kappa_B(\widetilde\rho_\bullet)\}
\,r_\bullet/\lVert r_\bullet\rVert_2,
& \lVert r_\bullet\rVert_2>0,\\[3pt]
0, & \lVert r_\bullet\rVert_2=0.
\end{cases}
\]

The exact clipping constants and zero-resultant tolerance are given in the
Supplementary Material. Define
$E_0^{\mathrm{com}}=\mathbf 1_K(\eta^{\mathrm{common}})^\top$ and
$G_0=\nabla f_{\tau^0}(E_0^{\mathrm{com}})$, where $f_{\tau^0}$ denotes the
smooth term $f_t$ in Eq.~\eqref{eq:fixed-responsibility-objective} evaluated
at the responsibilities $\tau^0$. A Karush--Kuhn--Tucker (KKT)-motivated
upper path scale evaluated at this approximate pooled fit is

\begin{equation}
\label{eq:kkt-endpoint}
\widetilde\lambda_{\max,\eta}^{(a)}
=1.05\max_{1\leq j\leq d}
\frac{\lVert (H_KG_0)_{\cdot j}\rVert_2}{w_j^{(a)}}.
\end{equation}

Conditional on $E_0^{\mathrm{com}}$, the uninflated maximum in
Eq.~\eqref{eq:kkt-endpoint} is the group threshold required by the
centered-block KKT condition. Because the pooled concentration is the capped Banerjee
approximation rather than the numerical solution of the pooled mean-resultant
equation, however, Eq.~\eqref{eq:kkt-endpoint} is not the exact KKT threshold
of the fully optimized all-common subproblem. We use it only as a
KKT-motivated path scale. An active cap additionally introduces a normal-cone
term, and the scale is not a global threshold for the observed mixture
objective. The path is

\begin{equation}
\label{eq:lambda-path}
\Lambda^{(a)}
=\{0\}\cup
\left\{
\widetilde\lambda_{\max,\eta}^{(a)}
10^{-4+4q/(L-2)}:
q=0,\ldots,L-2
\right\},
\end{equation}

with $L=240$. If $\widetilde\lambda_{\max,\eta}^{(a)}>0$, the 239 positive
values range from
$10^{-4}\widetilde\lambda_{\max,\eta}^{(a)}$ to
$\widetilde\lambda_{\max,\eta}^{(a)}$ on a logarithmic grid. E-CGL and
E-ACGL use the same relative range and grid size; their absolute penalty
values differ because the path scale contains the fixed method-specific
weights. If the endpoint is zero, the implementation sets
$\Lambda^{(a)}=\{0\}$ and proceeds directly to candidate refitting. Warm starts
are used in increasing penalty order.

Let

\[
\mathcal I_{\mathrm{record}}^{(a)}
=\left\{\ell:
\widehat\Theta_{\ell}^{(a)}\text{ is finite, free of numerical failure, and cap-feasible}
\right\}.
\]

A fit at its iteration limit is logged as unconverged but is retained in the
provisional candidate family; a numerical failure is logged and skipped.
Three consecutive numerical failures terminate the path.
The path also stops after storing a support of size at most one, while the
all-common support is always included. This is a computational heuristic, not
a nested-path claim, so nonempty selection is conditional on the finite
recorded path. In a targeted full-path audit reported in the Supplement,
continuing both E-series methods in six representative conditions (12 paths
in total) through all remaining tuning values produced neither a new support
nor re-entry above one coordinate.

\subsubsection{Target-preserving support-constrained refit}

The penalized path identifies candidate supports. For a candidate support $S$,
the support-constrained refit omits the sparsity penalty, retains the finite
concentration cap, and uses the parameterization

\begin{equation}
\label{eq:refit-parameterization}
\eta_{kj}=b_j+c_{kj},
\qquad
\sum_{k=1}^K c_{kj}=0,
\qquad
c_{kj}=0\quad(j\notin S).
\end{equation}

Hence $j\notin S$ implies
$\eta_{1j}=\cdots=\eta_{Kj}=b_j$: an inactive posterior-score contrast coordinate
retains its estimated common baseline in the natural-parameter space.

The two stages are separated because the group penalty remains active on
nonzero contrast groups. For any exact cap-inactive stationary point
$\Theta_{\ell}^{(a),\circ}$ of the penalized objective at
$\lambda_{\eta,\ell}>0$, with
$\lVert H_KE_{\ell,\cdot j}^{(a),\circ}\rVert_2>0$, the contrast-space score
satisfies

\begin{equation}
\label{eq:active-penalized-score}
H_K\nabla_{E_{\cdot j}}\ell
\left(\Theta_{\ell}^{(a),\circ}\right)
=\lambda_{\eta,\ell}w_j^{(a)}
\frac{H_KE_{\ell,\cdot j}^{(a),\circ}}
{\lVert H_KE_{\ell,\cdot j}^{(a),\circ}\rVert_2}.
\end{equation}

Thus the penalized stationary point does not satisfy the active-contrast score
equation of the corresponding support-constrained likelihood. This equality
displays the sparsity-penalty term omitted during refitting. The two-stage
construction has a Gaussian-mixture precedent in Lasso--MLE
\citep{meynet2012thesis,meynet2012sparse} and regression analogies in the
relaxed Lasso and least-squares post-Lasso
\citep{meinshausen2007relaxed,belloni2013postlasso}. These precedents motivate
but do not theoretically validate the present refit. The vMF refit carries no
claim of finite-sample bias reduction or risk improvement, valid post-selection
inference, or global optimality.

For each active coordinate, the first $K-1$ contrasts are free and the last is
set to their negative sum. Conditional on the responsibilities, the
natural-parameter update targets maximization, over $b\in\mathbb R^d$ and the
free active contrasts, of

\begin{equation}
\label{eq:refit-q}
Q_S(\eta)
=\sum_{k=1}^K
\left[
\eta_k^\top r_k
+N_k\log C_d(\lVert\eta_k\rVert_2)
\right]
\end{equation}

The path fit is projected onto the support constraint and, when needed, scaled
to the concentration cap. The smooth constrained M-step uses the
limited-memory Broyden--Fletcher--Goldfarb--Shanno algorithm with bounds
(L-BFGS-B) \citep{byrd1995lbfgsb} and analytic gradients, followed by step
halving until the row-norm cap and observed-likelihood guard hold. This
quasi-Newton routine is used only for refitting; the penalized E-series update
remains proximal. The implementation-specific refit limits and stopping
tolerances are reported under ``E-series numerical controls'' in the
Supplementary Material.

A refit is eligible for BIC comparison only if it is finite and converged,
with no numerical failure; its inner optimizer reports success; it satisfies the
centered-support constraint to numerical tolerance; and it has
$\max_k\widehat\kappa_k<0.99\kappa_{\max}$. The explicit empty-support refit is
initialized from the recorded fit with the smallest support; ties use the
largest observed log-likelihood and then the fixed path order.

The resulting numerically converged support-constrained refit retains the
common baseline; no global
maximum of the nonconvex mixture likelihood is claimed.

\subsubsection{Support selection by BIC after refit}
\label{subsec:bic-refit}

Following the BIC principle of \citet{schwarz1978bic}, let $m_S=|S|$.
Conditional on the path-generated candidate collection, scoring the refitted
rather than penalized likelihood removes the direct effect of
$\lambda_\eta$-dependent shrinkage from the parameter estimates used in the
BIC criterion. The practical E-series nominal-dimension approximation is

\begin{equation}
\label{eq:nominal-df}
\operatorname{df}_{E}(S)
=d+(K-1)m_S+(K-1)\mathbf 1(m_S>0).
\end{equation}

The three terms count the coordinatewise common baselines, the active centered
contrasts, and the mixing proportions when nonzero component contrasts are
present. When $S=\varnothing$, all component densities coincide and the mixing
proportions are not identifiable from the observed density, so they are not
counted. On a generic identifiable stratum, this is a nominal parameter
dimension. We use it as a practical within-method BIC criterion,
not as an exact Schwarz marginal-likelihood approximation or an effective
degrees-of-freedom result. The candidate collection is data-dependent, the
constrained refits are local numerical solutions, and mixture strata can be
singular at component mergers, boundary mixing weights, or the all-common
model \citep{drton2017singularbic}.

For each recorded path point, define

\begin{equation}
\label{eq:path-support-candidates}
\begin{aligned}
S_\ell^{(a)}
&=\left\{
j:\lVert H_K\widehat E_{\ell,\cdot j}^{(a)}\rVert_2
>\varepsilon_{\mathrm{supp}}
\right\},\\
\mathcal S_{\mathrm{path}}^{(a)}
&=\operatorname{ordered\,unique}
\left\{S_\ell^{(a)}:
\ell\in\mathcal I_{\mathrm{record}}^{(a)}\right\},\\
\mathcal S_{\mathrm{cand}}^{(a)}
&=\mathcal S_{\mathrm{path}}^{(a)}\cup\{\varnothing\},\\
\mathcal S_{\mathrm{elig}}^{(a)}
&=\left\{S\in\mathcal S_{\mathrm{cand}}^{(a)}:
\widehat\Theta_S^{\mathrm{refit}}
\text{ passes all numerical checks}\right\}.
\end{aligned}
\end{equation}

Candidate construction, refit scoring, and final selection are then

\begin{equation}
\label{eq:refit-bic}
\begin{aligned}
\operatorname{BIC}^{\mathrm{refit}}(S)
&=-2\ell(\widehat\Theta_S^{\mathrm{refit}})
+\log(n)\operatorname{df}_{E}(S),\\
\widehat S^{(a)}
&=\arg\min_{S\in\mathcal S_{\mathrm{elig}}^{(a)}}
\operatorname{BIC}^{\mathrm{refit}}(S).
\end{aligned}
\end{equation}

The all-common candidate is always refitted and enters the comparison if
eligible. Nonempty candidates remain conditional on the finite path. For a
nonempty population support, the event
\[
S_\eta^\star\in\mathcal S_{\mathrm{path}}^{(a)}
\]
defines path coverage. Conditional BIC success evaluates exact recovery given
this event, whereas overall exact recovery does not condition on it. For
$S_\eta^\star=\varnothing$, coverage equals one by construction and exact
recovery is $\widehat S^{(a)}=\varnothing$.

When several path fits generate the same support, the one with the largest
observed log-likelihood is its penalized source and initializes the constrained
refit. If the source of the provisional BIC winner is unconverged, the complete
path is rebuilt with the extended allowance described above, unconverged
sources are excluded, and refit-BIC selection is repeated. Exact BIC ties use
the first eligible candidate in the fixed order. Neither rule uses class
labels.

The tuning parameter is therefore treated only as a candidate-generation
index. BIC after refit selects a distinct support and does not define or
require a uniquely interpreted estimate of $\lambda_\eta$.

BIC comparisons follow the within-method convention in
Section~\ref{subsec:bic-refit}; cross-method values are descriptive.

EBIC \citep{chen2008ebic} with $\gamma_{\mathrm{EBIC}}\in\{0.25,0.5,1\}$ and
alternative degrees-of-freedom rules are sensitivity analyses; they do not
remove mixture singularity, local optimization, or data-dependent candidate
construction.

Algorithm~\ref{alg:ecgl-path} writes out the E-step, the guarded proximal
generalized M-step, path construction, and refit-based support selection.
E-CGL is obtained with unit weights; E-ACGL uses the fixed dense-pilot weights
in \eqref{eq:eacgl-penalty}. All symbols and updates in the listing are defined
in the preceding subsections.
The fixed-$\lambda$ subroutine returns $\mathtt{FAIL}$ only after a numerical
guard or finite-value check fails. Reaching an iteration limit returns a fit
that has not failed numerically but has its convergence flag set to false, so
that failure and nonconvergence remain distinct; Step~3 retains the latter
provisionally and Step~4 validates the selected source.
Lines 4--17 of the listing define the fixed-$\lambda_\eta$ subroutine
$\operatorname{Fit}_{\lambda_\eta}^{(a)}$ used in the subsequent path steps.

\begin{boxedalgorithm}{Guarded proximal GEM and support selection}
\label{alg:ecgl-path}
\textbf{Require:} row-normalized $X\in\mathbb R^{n\times d}$; $K$;
$a\in\mathcal A$; dense starts; $L=240$; $\kappa_{\max}$; iteration limits
$T_0=100$, $T_{\mathrm{retry}}=440$, and $U_{\max}$; the stopping tolerances in
Eqs.~\eqref{eq:pg-stopping}--\eqref{eq:gem-stopping};
$\varepsilon_{\mathrm{maj}}$, $\varepsilon_{\mathrm{acc}}$,
$\varepsilon_{\mathrm{supp}}$; and the backtracking limit.\par
\textbf{Ensure:} $\widehat S^{(a)}$ and
$\widehat\Theta_{\widehat S^{(a)}}^{\mathrm{refit}}$.
\par\smallskip
{\renewcommand{\arraystretch}{0.96}
\begin{tabularx}{\textwidth}{@{}r@{\hspace{0.60em}}X@{}}
1: & \textbf{Step 1: Initialization.}
\tabularnewline
2: & \textbf{(a) Dense pilot:} select $\widehat\Theta^{\mathrm D}$ by the
multistart continuation-and-fallback rule above.
\tabularnewline
3: & \textbf{(b) Fixed weights:} compute and fix
$w_j^{(a)}$, $j=1,\ldots,d$, according to
Eq.~\eqref{eq:eacgl-penalty}.
\tabularnewline
4: & \textbf{Step 2: Fixed-$\lambda_\eta$ proximal GEM.}
\tabularnewline
5: & Initialize $\Theta^{(0)}$ from the supplied warm start and set
$\mathtt{converged}=\mathtt{FALSE}$.
\tabularnewline
6: & $\mathbf{for}\ t=0,\ldots,T_0-1\ \mathbf{do}$
\tabularnewline
7: & \quad\textbf{Step 2.1: E-step.} Compute $\tau_{ik}^{(t)}$,
$N_k^{(t)}$, and $r_k^{(t)}$ using Eq.~\eqref{eq:responsibilities} and the
definitions immediately below it.
\tabularnewline
8: & \quad\textbf{Step 2.2: Update mixing proportions.}
Compute $\pi^{\mathrm{cand}}$ by Eq.~\eqref{eq:mixing-update}.
\tabularnewline
9: & \quad\textbf{Step 2.3: Update natural parameters.}
Initialize the inner iteration at the current natural-parameter matrix.
\tabularnewline
10: & $\displaystyle
\qquad\mathbf{for}\ u=0,\ldots,U_{\max}-1\ \mathbf{do}$
\tabularnewline
11: & $\displaystyle
\qquad\quad s\leftarrow(\max_kN_k^{(t)})^{-1},\qquad
V\leftarrow E^{(u)}-s\nabla f_t(E^{(u)})$.
\tabularnewline
12: & \qquad Apply the centered coordinatewise proximal update in
Eq.~\eqref{eq:centered-prox-update} to $V$, for $j=1,\ldots,d$.
\tabularnewline
13: & Halve $s$ and recompute $V$ and $E^{(u+1)}$ until the majorization and
concentration-cap conditions hold; if no admissible proposal is found,
$\mathbf{return}\ \mathtt{FAIL}$.
\tabularnewline
14: & If Eq.~\eqref{eq:pg-stopping} holds, $\mathbf{break}$;
$\mathbf{end\ for}$; set
$\Theta^+\leftarrow(\pi^{\mathrm{cand}},E^{(u+1)})$.
\tabularnewline
15: & \quad\textbf{Step 2.4: Acceptance and convergence.}
Accept $\Theta^+$ only if Eq.~\eqref{eq:outer-acceptance} holds; otherwise
$\mathbf{return}\ \mathtt{FAIL}$.
\tabularnewline
16: & Set $\Theta^{(t+1)}\leftarrow\Theta^+$; if
Eq.~\eqref{eq:gem-stopping} holds, set
$\mathtt{converged}=\mathtt{TRUE}$ and $\mathbf{break}$.
\tabularnewline
17: & $\mathbf{end\ for}$; at the iteration limit, retain the current finite
fit with $\mathtt{converged}=\mathtt{FALSE}$.
\tabularnewline
18: & \textbf{Step 3: Increasing-penalty path.}
\tabularnewline
19: & $\displaystyle
\widehat\Theta_0^{(a)}\leftarrow
\operatorname{Fit}_{\lambda_\eta=0}^{(a)}
(\widehat\Theta^{\mathrm D};w^{(a)});\qquad
\widehat\Theta_0^{(a)}=\mathtt{FAIL}\Rightarrow
\mathbf{return}\ \mathtt{FAIL}$.
\tabularnewline
20: & Set $\mathcal I_{\mathrm{record}}^{(a)}\leftarrow\{0\}$ and form
$S_0^{(a)}$ from $\widehat\Theta_0^{(a)}$ by
Eq.~\eqref{eq:path-support-candidates}.
\tabularnewline
21: & Construct the KKT-motivated scale
$\widetilde\lambda_{\max,\eta}^{(a)}$ and the single path
$\Lambda^{(a)}$ from Eqs.~\eqref{eq:kkt-endpoint}--\eqref{eq:lambda-path}.
\tabularnewline
22: & $\mathbf{for}$ each positive $\lambda_{\eta,\ell}\in\Lambda^{(a)}$ in
increasing order $\mathbf{do}$ initialize Step~2 from the preceding accepted fit.
\tabularnewline
23: & Repeat Step~2. If the fit is finite, nonfailed, and cap-feasible, include
$\ell$ in $\mathcal I_{\mathrm{record}}^{(a)}$ and form $S_\ell^{(a)}$ by
Eq.~\eqref{eq:path-support-candidates}.
\tabularnewline
24: & Log an iteration-limit fit as unconverged. On numerical failure, retain
the last accepted warm start and continue; three consecutive failures terminate the path. After recording
$|S_\ell^{(a)}|\leq1$, stop the path. $\mathbf{end\ for}$
\tabularnewline
25: & \textbf{Step 4: Target-preserving refit and support selection.}
\tabularnewline
26: & Construct $\mathcal S_{\mathrm{cand}}^{(a)}$ by
Eq.~\eqref{eq:path-support-candidates}; for each candidate, compute
$\widehat\Theta_S^{\mathrm{refit}}$ under
Eq.~\eqref{eq:refit-parameterization}, apply the stated numerical checks, and
then form $\mathcal S_{\mathrm{elig}}^{(a)}$.
\tabularnewline
27: & Select $\widehat S^{(a)}$ by Eq.~\eqref{eq:refit-bic}; if
$\mathcal S_{\mathrm{elig}}^{(a)}=\varnothing$, $\mathbf{return}\ \mathtt{FAIL}$.
\tabularnewline
28: & If the penalized source of $\widehat S^{(a)}$ is unconverged, rebuild
Step~3 with $T_0+T_{\mathrm{retry}}$, retain converged path sources only, and
repeat Steps~4 and 27.
\tabularnewline
29: & $\displaystyle
\mathbf{return}\ (\widehat S^{(a)},
\widehat\Theta_{\widehat S^{(a)}}^{\mathrm{refit}})$.
\tabularnewline
\end{tabularx}}
\end{boxedalgorithm}
In the proximal update, the centered term is zero when
$H_KV_{\cdot j}=0$. The guard logic is common to both methods, whereas the
penalized criteria evaluated by the guards are method-specific through
$w_j^{(a)}$. Refit BIC compares eligible supports only within a method.

\subsubsection{Selection of the number of components}
\label{subsec:k-selection}

Component-number selection is not part of the E-CGL support estimator. The
simulations therefore fix $K$ at its generating value. When $K$ is unknown,
we first fit dense multistart mixtures over a prespecified set $\mathcal K$ and
choose a resolution using a predeclared label-free density or stability
criterion oriented so that smaller values are preferred. A stability score is
therefore entered through its negative or another decreasing transformation.
E-CGL support selection is then performed at that resolution:

\[
\widehat K
=\arg\min_{K\in\mathcal K}\mathcal C_{\mathrm{dense}}(K),
\qquad
\widehat S_{\widehat K}
=\arg\min_{S\in\mathcal S_{\mathrm{elig}}^{(\mathrm{CGL})}(\widehat K)}
\operatorname{BIC}^{\mathrm{refit}}(\widehat K,S).
\]

This two-stage diagnostic has no component-number consistency claim; competing
resolutions should be reported when density fit and stability disagree.

\subsubsection{Computational implementation}

The implementation uses log-sum-exp stabilization in the E-step, warm starts
along the penalty path, and explicit convergence and constraint checks in the
centered-support refit. Optional compiled helpers use the Rcpp interface
\citep{eddelbuettel2011rcpp} for low-level numerical operations without
changing the statistical algorithm. R-only and compiled-helper outputs were
checked to numerical tolerance before runtime comparisons.

The current computational claims are limited to empirical equality and runtime
diagnostics. The method does not rely on an Rcpp implementation for its
statistical definition.

The finite parameter space $\mathcal T_{\kappa_{\max}}$ is used in dense
initialization, the penalized path, and support-constrained refitting.
Backtracking rejects any proposed E-series iterate outside this space.
A converged refit that reaches 99\% of the bound is retained for diagnosis but
is not eligible for information-criterion selection. This guard defines the
finite computational problem; it is not an additional sparsity penalty.

\subsection{Structural properties}
\label{sec:structural}

The following statements justify the centered support estimand for a fixed
number of components. They do not imply uniqueness of a fitted mixture, which
remains subject to finite-mixture identifiability and label permutation. Full
proofs are given in \hyperref[app:proofs]{Appendix A}.

\begin{proposition}[Centered decomposition and pairwise support]
\label{prop:centered-support}
For any $v\in\mathbb R^K$, there is a unique decomposition

\[
v=\bar v\mathbf1_K+c,
\qquad
\bar v=K^{-1}\mathbf1_K^\top v,
\qquad
\mathbf1_K^\top c=0.
\]

For $v=\eta_{\cdot j}$, the following statements are equivalent:

\[
c=0;
\qquad
\eta_{1j}=\cdots=\eta_{Kj};
\qquad
\eta_{kj}-\eta_{\ell j}=0\quad\text{for all }k<\ell.
\]

Moreover,

\[
\lVert c\rVert_2^2
=\frac1K\sum_{k<\ell}(\eta_{kj}-\eta_{\ell j})^2.
\]
Consequently,

\[
S_{\eta}
=\bigcup_{k<\ell}\{j:\eta_{kj}-\eta_{\ell j}\neq0\},
\]

the union of the feature-dependent linear supports in all pairwise posterior
log-odds.
\end{proposition}

Proposition~\ref{prop:centered-support} clarifies the scope of the estimand. If
$c_{\cdot j}=0$, the coefficient of $x_j$ cancels from every pairwise
posterior log-odds. The common baseline is nevertheless retained in $\eta_k$
and may affect the intercept through $\lVert\eta_k\rVert_2$; it is therefore
not irrelevant to the component density.

\begin{proposition}[Natural-parameter and directional decomposition]
\label{prop:canonical-directional}
Let $\kappa=(\kappa_1,\ldots,\kappa_K)^\top$,
$\bar\kappa=K^{-1}\mathbf1_K^\top\kappa$, and
$h=\kappa-\bar\kappa\mathbf1_K$. Write
\[
\bar\mu=K^{-1}M^\top\mathbf1_K,
\qquad
U=H_KM.
\]
Then $M=\mathbf1_K\bar\mu^\top+U$ and
\[
H_KE
=\bar\kappa U+h\bar\mu^\top+H_K\operatorname{diag}(h)U.
\]
The three terms represent directional heterogeneity, concentration
heterogeneity along the mean directional loading, and their interaction.
For support interpretations involving $M$ and $S_\mu$, assume
$\kappa_k>0$ for every $k$.
\end{proposition}

If $\kappa_1=\cdots=\kappa_K=\kappa>0$, then
$H_KE=\kappa H_KM$ and $S_\eta=S_\mu$. If
$\mu_1=\cdots=\mu_K=\bar\mu$ and $h\neq0$, then
$S_\mu=\varnothing$ while $S_\eta=\operatorname{supp}(\bar\mu)$. Conversely,
heterogeneous concentrations can make $H_KM_{\cdot j}\neq0$ even when
$H_KE_{\cdot j}=0$. In the general decomposition the three terms may cancel,
so $S_\eta$ is not identified with the union of their separate supports.

\begin{proposition}[Label equivariance and support invariance]
\label{prop:label-invariance}
Let $P$ be a $K\times K$ permutation matrix. Then

\[
H_KP=PH_K,
\qquad
\lVert H_KP\eta_{\cdot j}\rVert_2
=\lVert H_K\eta_{\cdot j}\rVert_2.
\]

The penalties $\mathcal P_a$, $a\in\mathcal A$, are invariant to component
relabeling, their proximal maps are equivariant, and $S_{\eta}$ is unchanged.
For E-ACGL, the dense pilot is relabeled consistently before its fixed weights
are computed.
\end{proposition}

For a nonzero natural parameter, the component representation is recovered as
$\kappa_k=\lVert\eta_k\rVert_2$ and
$\mu_k=\eta_k/\lVert\eta_k\rVert_2$. At $\eta_k=0$, the vMF component is
uniform and its direction is not identified. These algebraic statements do
not establish finite-mixture identifiability or uniqueness of the fitted
mixture. Component-specific interpretation assumes distinct component
distributions and positive mixing weights and is made only up to label
permutation. If two components have identical natural parameters, their
individual mixing weights are not separately identified. Support-recovery
statements are conditional on fixed $K$.

\subsection{Algorithmic properties}

\begin{proposition}[Closed-form proximal update]
\label{prop:centered-prox}
Let $J_K=K^{-1}\mathbf1_K\mathbf1_K^\top$ and $H_K=I_K-J_K$. For
$\xi\geq0$ and $v\in\mathbb R^K$,

\[
\operatorname{prox}_{\xi\lVert H_K\cdot\rVert_2}(v)
=J_Kv+
\left(1-\frac{\xi}{\lVert H_Kv\rVert_2}\right)_+H_Kv,
\]

where the second term is defined as zero when $H_Kv=0$.
\end{proposition}

The common baseline is unchanged. Group soft-thresholding is applied only
to the centered contrast, matching the coordinate-level definition of
$S_{\eta}$.

\begin{proposition}[Smooth majorization and proximal descent]
\label{prop:majorization}
Fix responsibilities with $N_k>0$. The function
$f_t$ in Section~\ref{sec:fixed-responsibility} is convex and has a Lipschitz
gradient with constant $L_t\leq\max_kN_k$. Hence, for $s\leq1/L_t$, the
proximal-gradient update in Section~\ref{sec:centered-pg-update} does not
increase $F_t^{(a)}$ in exact arithmetic. More generally, an update accepted with
majorization tolerance $\varepsilon_{\mathrm{maj}}$ satisfies

\[
F_t^{(a)}(E^{(u+1)})
\leq F_t^{(a)}(E^{(u)})+\varepsilon_{\mathrm{maj}}.
\]
\end{proposition}

The backtracking rule starts at $s=1/\max_kN_k$ and halves the step until this
majorization condition holds numerically.

\begin{theorem}[Fixed-responsibility proximal convergence]
\label{thm:fixed-responsibility}
Let $d\geq2$. Fix responsibilities such that $N_k>0$ and
$\lVert r_k\rVert_2<N_k$ for every $k$, fix $a\in\mathcal A$ with finite
positive weights, and fix a finite $\lambda_\eta\geq0$. Then
$F_t^{(a)}$ on $\mathbb R^{K\times d}$ is strictly convex and coercive and therefore
has a unique unconstrained minimizer $E_t^\star$. If $E_t^\star$ lies strictly
inside the concentration bound, it is also the unique solution of the bounded
subproblem. The unconstrained proximal-gradient method with the exact
majorization backtracking rule ($\varepsilon_{\mathrm{maj}}=0$) converges to
$E_t^\star$,
and its objective gap
is $O(u^{-1})$ after $u$ inner iterations.
\end{theorem}

Strict convexity follows from the positive-definite vMF covariance at finite
natural parameters. The condition $\lVert r_k\rVert_2<N_k$ makes the radial
linear term dominate the logarithmic Bessel term, yielding coercivity; the
global Hessian bound then gives the stated proximal-gradient result. The
complete argument is given in \hyperref[app:proofs]{Appendix A}.

The theorem concerns the limiting fixed-responsibility problem with an
inactive cap. Finite inner tolerances are inexact, and an active cap would
require projected-proximal analysis \citep{beck2009fista,parikh2014proximal}.

\begin{proposition}[Penalized GEM monotonicity]
\label{prop:gem-monotonicity}
Let $\tau^{(t)}$ be the posterior responsibilities under $\Theta^{(t)}$. If a
candidate $\Theta^{(t+1)}$ satisfies

\[
Q_{\lambda_\eta}^{(a)}(\Theta^{(t+1)}\mid\tau^{(t)})
\geq
Q_{\lambda_\eta}^{(a)}(\Theta^{(t)}\mid\tau^{(t)}),
\]

then, in exact arithmetic,

\[
\mathcal L_{\lambda_\eta}^{(a)}(\Theta^{(t+1)})
\geq
\mathcal L_{\lambda_\eta}^{(a)}(\Theta^{(t)}).
\]
\end{proposition}

The implementation permits decreases of at most $10^{-8}$, so finite traces
are only numerically near-monotone. Exact monotonicity alone does not imply
parameter convergence or global optimality; the next theorem instead gives a
conditional stationarity result.

For a responsibility array $\tau$, let $f_\tau$ denote the smooth
fixed-responsibility objective in Section~\ref{sec:fixed-responsibility}. Define

\begin{equation}
\label{eq:block-residual-definitions}
\begin{aligned}
\mathcal G_s(E;\tau)
&=\frac{1}{s}
\left[
E-\operatorname{prox}_{s\lambda_\eta\mathcal P_a}
\{E-s\nabla f_\tau(E)\}
\right],\\
\mathcal R_\pi(\pi;\tau)
&=\pi-\frac{N(\tau)}{n},
\qquad
N_k(\tau)=\sum_{i=1}^n\tau_{ik}.
\end{aligned}
\end{equation}

For Theorem~\ref{thm:idealized-gem}, $s_t$ is the diagnostic step size used to
evaluate the natural-parameter residual at the final accepted inner iterate;
it need not equal every trial step used during backtracking.

\begin{theorem}[Conditional stationarity under vanishing block residuals]
\label{thm:idealized-gem}
Fix $a\in\mathcal A$ with finite positive weights and fix a finite
$\lambda_\eta\geq0$.
Consider an idealized GEM sequence
$\{\Theta^{(t)}\}$ for which

\begin{equation}
\label{eq:stationarity-assumptions}
\begin{gathered}
\tau^{(t)}=\tau(\Theta^{(t)}),
\qquad
\Theta^{(t)},\Theta^{(t+1)}
\in\mathcal C\Subset\operatorname{ri}(\mathcal T_{\kappa_{\max}}),\\
\inf_{t,k}N_k(\tau^{(t)})>0,
\qquad
0<\underline s\leq s_t\leq\overline s<\infty,\\
\lVert\Theta^{(t+1)}-\Theta^{(t)}\rVert\longrightarrow0,
\end{gathered}
\end{equation}

where $\operatorname{ri}$ denotes relative interior, the norm on $\Theta$ is
the Euclidean product norm, and
$\mathcal C\Subset\operatorname{ri}(\mathcal T_{\kappa_{\max}})$ means that
$\mathcal C$ is a compact subset of that relative interior. Suppose further
that the block residuals satisfy

\begin{equation}
\label{eq:vanishing-block-residuals}
\lVert\mathcal G_{s_t}(E^{(t+1)};\tau^{(t)})\rVert_F
+\lVert\mathcal R_\pi(\pi^{(t+1)};\tau^{(t)})\rVert_2
\longrightarrow0.
\end{equation}

Then every accumulation point $\bar\Theta=(\bar\pi,\bar E)$ satisfies the
interior KKT conditions

\begin{equation}
\label{eq:kkt-stationarity}
\begin{gathered}
0\in-\nabla_E\ell(\bar\Theta)
+\lambda_\eta\partial\mathcal P_a(\bar E),\\
\nabla_\pi\ell(\bar\Theta)
=\nu\mathbf1_K,
\qquad
\mathbf1_K^\top\bar\pi=1,
\qquad \text{for some }\nu\in\mathbb R.
\end{gathered}
\end{equation}

Thus every accumulation point is stationary for the finite-parameter
penalized observed objective. The stationarity statement is equivariant under
component-label permutation.
\end{theorem}

The proof passes to a convergent subsequence, uses continuity of the
responsibility and proximal maps, and applies Fisher's identity to obtain the
observed-score KKT system. Details are given in
\hyperref[app:proofs]{Appendix A}.

This conditional local result is not a convergence claim for every
finite-tolerance path. Active bounds, failed path points, nonvanishing
residuals, or component degeneracy fall outside its assumptions; see
\citep{dempster1977em,green1990penalizedem,wu1983em,tseng2004emproximal}
for related EM convergence theory.

\subsection{Computational complexity}

For dense $X$, one E-step costs $O(nKd)$; for a sparse matrix it costs
$O\{K\operatorname{nnz}(X)\}$, where $\operatorname{nnz}(X)$ is the number of
nonzero entries of $X$. Forming the sufficient statistics has the same
order, and one centered proximal-gradient update costs $O(Kd)$. Excluding the
input matrix, the main working memory is $O(nK+Kd)$, in addition to
$O(nd)$ for dense $X$ or $O\{\operatorname{nnz}(X)\}$ for sparse $X$. If the
single path evaluates $L$ tuning points, let $T_{\mathrm{GEM}}$ denote the
average number of outer GEM iterations per path point and
$T_{\mathrm{PG}}$ the average number of inner proximal-gradient updates per
GEM iteration. Its leading cost is
\[
O\left[LT_{\mathrm{GEM}}
\left\{K\operatorname{nnz}(X)+T_{\mathrm{PG}}Kd\right\}\right]
\]
for sparse input, plus the costs of distinct-support refits. These are
operation counts rather than runtime superiority claims.
Both E-CGL and E-ACGL use $L=240$ before early stopping. Multiple dense
starts add their own E-step and M-step costs before the path is constructed.

\section{Numerical studies}
\label{sec:simulations}

\subsection{Design and evaluation criteria}

The primary experiment evaluates recovery of posterior-score contrast support
under varying overlap and concentration heterogeneity. Observations are drawn
from a balanced $K=4$ vMF mixture in $d=200$ dimensions. The coordinates are
partitioned as
\[
\{1,\ldots,d\}=S_C\cup S_D\cup S_N,
\qquad
(q_C,q_D,q_N)=(4,16,180),
\]
where $S_C$ contains common natural-parameter coefficients, $S_D$ contains
centered component contrasts, and $S_N$ contains globally null coordinates.
Let $q_{\mathrm{block}}=4$ be the number of contrast coordinates assigned to
each component block. For a contrast coordinate $j$ in block $g$, define
\[
u_{kj}=
\begin{cases}
1, & k=g,\\
-1/(K-1), & k\ne g.
\end{cases}
\]
and set $u_k=(u_{kj}:j\in S_D)\in\mathbb R^{q_D}$. For a candidate
common-background norm satisfying
$0\leq A<\min_{1\leq k\leq K}\kappa_k$, the natural parameters are generated by
\[
\eta_{kj}(A)=
\begin{cases}
A/\sqrt{q_C}, & j\in S_C,\\
\alpha_k(A)u_{kj}, & j\in S_D,\\
0, & j\in S_N,
\end{cases}
\qquad
\alpha_k(A)=\frac{\sqrt{\kappa_k^2-A^2}}{\lVert u_k\rVert_2}.
\]
Hence $\lVert\eta_k(A)\rVert_2=\kappa_k$ and
$\mu_k(A)=\eta_k(A)/\kappa_k$. The clustering difficulty is
calibrated by the oracle Bayes error
\[
e_B=\Pr_{\Theta^\star}
\left\{
\arg\max_k\Pr(Z=k\mid X;\Theta^\star)\ne Z
\right\}.
\]
We use $e_B\in\{0.025,0.05,0.10\}$, $n\in\{300,1000\}$, and either
\[
\kappa=(45,45,45,45)
\quad\text{or}\quad
\kappa=(30,40,50,60).
\]
The common-to-contrast allocation is calibrated by Monte Carlo bisection while
the prescribed concentrations and support sets are held fixed. Independent
validation samples of size 200,000 gave achieved errors between 0.02566 and
0.10035; the largest absolute discrepancy from its target was 0.00172.

Each of the 12 cells contains 100 paired repetitions and an independent test
sample of size 5,000. The common comparison set consists of spherical
$k$-means, dense vMF mixtures with shared and component-specific
concentrations, the Rossi-style entrywise mean-direction lasso (M-L), E-CGL,
and E-ACGL. The centered
mean-direction group lasso (M-CGL) is included as a
directional companion in the four prespecified
$e_B=0.05$ cells. E-CGL and E-ACGL use the single 240-point,
zero-augmented KKT-motivated path in \eqref{eq:lambda-path}. M-CGL uses a separate
240-point dense-score-proxy path from
$10^{-4}\lambda_{\mathrm{proxy},\mu}$ to
$1.25\lambda_{\mathrm{proxy},\mu}$, where the proxy scale is defined in
Supplementary Section S8. M-L retains its Rossi-style event-driven path, with
at most 240 path points, including the zero-penalty fit, in the simulation,
because its threshold events are not indexed by either of the centered-group
path scales.
All mixture fits use 10 dense starts and Rcpp helpers; sparse mixture fits use
BIC after target-specific support refitting.
Across the common 24-cell rerun---the 12 primary cells, three additional
estimand-separation cells, four $e_B=0.05$ sample-size cells, and five stress
cells---four of the 2,400 dataset-replicate dense
free-concentration pilots shared by the two E-series methods reached the
200-iteration limit. Continuing each terminal iterate to convergence left all eight
corresponding E-CGL and E-ACGL supports unchanged; the largest absolute BIC
change was 0.027.
An exhaustive audit of the selected penalized source in all 4,800 E-series
method--replicate fits identified 37 sources that had reached the original
100-iteration limit; 32 selected supports had no converged source on that
recorded path. We therefore reran only those 37 method--replicate paths with a
total limit of 540 outer iterations per point and admitted converged path
points only. All 37 corrected
sources and refits converged; 34 supports changed, including 32 changes in
support size. The primary 12-cell results were unaffected. All additional
and stress results reported below use the corrected archive; full audit
details are given in the Supplementary Material.

Table~\ref{tab:simulation-methods} summarizes the estimand and role of each
method. The common comparison set is evaluated in all 12 primary cells;
M-CGL is restricted to the four prespecified $e_B=0.05$ companion cells.

\FloatBarrier
\begin{table}[H]
\centering
\small
\caption{Methods, support targets, and roles in the numerical comparisons.}
\label{tab:simulation-methods}
\resizebox{\textwidth}{!}{%
\begin{tabular}{lll}
\toprule
Method & Support target & Role \\
\midrule
Spherical $k$-means & none & external clustering baseline \\
Dense vMF, shared $\kappa$ & none & concentration-restricted baseline \\
Dense vMF, free $\kappa_k$ & none & likelihood baseline \\
M-L & $S_P$ & published prototype-sparsity reference \\
M-CGL & $S_\mu$ & directional companion \\
E-CGL & $S_\eta$ & primary proposed estimator \\
E-ACGL & $S_\eta$ & secondary adaptive extension \\
\bottomrule
\end{tabular}}
\end{table}

For clarity, the two mean-direction comparators use different penalty
functionals. The Rossi-style entrywise prototype penalty is

\[
\mathcal P_{\mathrm{M-L}}(M)
=\sum_{k=1}^K\sum_{j=1}^d|\mu_{kj}|,
\qquad
\lVert\mu_k\rVert_2=1,
\]

whereas the directional companion uses

\[
\mathcal P_{\mathrm{M-CGL}}(M)
=\sum_{j=1}^d
\lVert c_{\cdot j}^{(\mu)}\rVert_2,
\qquad
\lVert\mu_k\rVert_2=1.
\]

The respective fitted criteria subtract
$\lambda_\mu\mathcal P_{\mathrm{M-L}}(M)$ and
$\lambda_\mu\mathcal P_{\mathrm{M-CGL}}(M)$ using their method-specific
paths.

In the primary simulations, M-L estimates component-specific concentrations so
that prototype selection is not coupled to a false shared-concentration
restriction. The Rossi bridge and real-data M-L fit instead use the published
shared-concentration specification; their results are conditional on that
restriction. M-CGL introduces an auxiliary matrix
$\Xi\in\mathbb R^{K\times d}$ subject to $\Xi=H_KM$, thereby separating a
product-of-spheres direction update from centered group thresholding and a
scaled-dual update before numerical concentration updates. At feasibility,
$\Xi_{\cdot j}=c_{\cdot j}^{(\mu)}$. Its target-preserving refit enforces
$c_{\cdot j}^{(\mu)}=0$ outside the selected
directional-contrast support. The augmented-Lagrangian updates and numerical residual
criteria are given in Supplementary Section S8, \emph{Directional companion
implementation}.

Clustering agreement is evaluated by the adjusted Rand index (ARI)
\citep{hubert1985comparing} and normalized mutual information (NMI)
\citep{strehl2002cluster}. The implementation uses geometric-mean entropy
normalization for NMI; larger values indicate closer agreement. Predictive fit
is evaluated by the per-observation test negative log-likelihood (NLL)
\[
\operatorname{NLL}_{\mathrm{test}}
=-\frac{1}{n_{\mathrm{test}}}
\sum_{i=1}^{n_{\mathrm{test}}}
\log f(X_i;\widehat\Theta),
\]
where $f$ is the vMF mixture density with respect to surface-area measure on
the unit sphere. Such a density can exceed one, so its reported NLL can be
negative. Let $\widehat S$ denote the coordinate set selected by the method
under evaluation. For the E-series methods, $\widehat S=\widehat S_\eta$;
for M-L and M-CGL, it estimates $S_P$ and $S_\mu$, respectively.
Posterior-score contrast support recovery is summarized
against $S_\eta^\star=S_D$, first by the selected-coordinate counts
\[
\widehat q_C=|\widehat S\cap S_C|,
\qquad
\widehat q_D=|\widehat S\cap S_D|,
\qquad
\widehat q_N=|\widehat S\cap S_N|,
\]
and then by the true-positive rate (TPR), false-positive rate (FPR), precision,
and $F_1$ score, together with the exact-support rate. Using
$(\eta^c)_{kj}=c_{kj}$ for a centered natural-parameter
entry, write
$c_{kj}^{\star}=\eta_{kj}^{\star}
-K^{-1}\sum_h\eta_{hj}^{\star}$ and
$\widehat c_{kj}=\widehat\eta_{kj}
-K^{-1}\sum_h\widehat\eta_{hj}$. Centered natural-parameter estimation is
summarized, after MSE-optimal component alignment, by
\[
\operatorname{MSE}_{\eta^c}
=\min_{P\in\mathfrak P_K}
\frac{1}{Kd}
\left\lVert P H_K\widehat E-H_KE^\star\right\rVert_F^2,
\]
where $\mathfrak P_K$ is the set of $K\times K$ permutation matrices. Selected
table entries are accompanied by Monte Carlo
standard errors (MCSEs), computed across the 100 repetitions in each cell. The
primary support tables evaluate every sparse
procedure against $S_\eta^\star$ to answer a common posterior-score contrast question. Those
entries are therefore cross-target diagnostics for M-L and M-CGL.
Target-specific recovery is reported separately: M-L is evaluated against
$S_P$ in the summary below, and M-CGL is evaluated against
$S_\mu$ in Section~\ref{sec:estimand-comparison}. Consequently, $F_{1,P}$,
$F_{1,\mu}$, and $F_{1,\eta}$ are not interchangeable method rankings.

\subsection{Posterior-score contrast support recovery}

Tables~\ref{tab:primary-n300} and \ref{tab:primary-n1000} report estimation,
posterior-score selection, and clustering over the 12-cell grid.
Figs.~\ref{fig:primary-ari}--\ref{fig:primary-mse} show the replication-level
ARI and $\log_{10}(\mathrm{MSE}_{\eta^c})$ distributions. Common-target support
metrics use $S_\eta^\star$; M-L and M-CGL are also assessed against their own
targets below.

Across 11 of the 12 primary cells, E-CGL selected between 16.02 and 16.67
coordinates and attained $F_{1,\eta}\ge 0.981$. The exception combined
$n=300$, $e_B=0.10$, and heterogeneous concentrations. In that cell E-CGL
selected 20.97 coordinates and attained $F_{1,\eta}=0.816$; all six
$n=1000$ cells had $F_{1,\eta}\ge0.988$. M-L retained 199.42--199.88 of 200
coordinates. Its low cross-target $F_{1,\eta}$ partly reflects a different
estimand, but its own-target $F_{1,P}=0.182$ also shows near-dense prototype
selection. ARI decreased with $e_B$ even where support recovery remained
accurate.
Under the common single-path rule, E-ACGL
selected 15.23--16.33 coordinates at $n=300$ and 16.01--16.06 at $n=1000$;
its corresponding $F_{1,\eta}$ ranges were 0.947--0.996 and 0.998--1.000.

\begin{figure}[t]
\centering
\includegraphics[width=\textwidth]{figure_01.pdf}
\caption{Clustering accuracy across the primary simulation grid. Boxplots show
ARI over 100 paired replications in each cell. Columns give the target oracle
Bayes error, and rows give sample size and concentration pattern. D-S and D-F
denote dense vMF mixtures with shared and component-specific concentrations,
respectively. M-CGL is omitted
because it was evaluated only in the four
prespecified $e_B=0.05$ cells; its target-specific results are reported in
Table~\ref{tab:estimand-separation}. Spherical $k$-means is omitted from the
figure for readability and is reported in Tables~\ref{tab:primary-n300} and
\ref{tab:primary-n1000}.}
\label{fig:primary-ari}
\end{figure}

\begin{figure}[t]
\centering
\includegraphics[width=\textwidth]{figure_02.pdf}
\caption{Centered natural-parameter estimation error across the primary simulation
grid. Boxplots show $\log_{10}(\mathrm{MSE}_{\eta^c})$ over 100 paired
replications in each cell, where $\eta^c=H_KE$. Columns give the target oracle
Bayes error, and rows give sample size and concentration pattern. D-S and D-F denote dense vMF
mixtures with shared and component-specific concentrations, respectively.
M-CGL is omitted because it was evaluated only in the four prespecified
$e_B=0.05$ cells; its target-specific results are reported in
Table~\ref{tab:estimand-separation}.}
\label{fig:primary-mse}
\end{figure}

\input{../simulations/tables/csda_main_section42_n300_260804.tex}

\input{../simulations/tables/csda_main_section42_n1000_260804.tex}

\FloatBarrier
\subsection{Directional versus posterior-score estimands}
\label{sec:estimand-comparison}

For the simulation truths in this subsection, let
$q_\mu^\star=|S_\mu|$ and $q_\eta^\star=|S_\eta|$.

Four diagnostics separate directional and posterior-score heterogeneity. Under
common concentration, $S_\mu=S_\eta$. Under pure concentration
heterogeneity, the component directions are common, so $S_\mu=\varnothing$
while $S_\eta$ contains 16 coordinates; this design uses
$\kappa=(10,30,80,200)$. The shared-natural-parameter design uses
$\kappa=(30,40,50,60)$ and contains
80 common natural-parameter coordinates and 20 posterior-score contrast coordinates.
The crossed design uses the same heterogeneous concentrations and contains
four $\eta$-only, four $\mu$-only, eight jointly active, and eight null
coordinates. Table~\ref{tab:estimand-separation}
compares recovery against both population targets in these four designs.

\FloatBarrier
\begin{table}[H]
\centering
\small
\caption{Directional and posterior-score contrast recovery in the estimand-separation
designs. Entries are means over 100 replications, and the two target-specific
$F_1$ scores are evaluated against their respective population targets. The exact-support column uses $S_\mu$ for
M-CGL and $S_\eta$ for E-CGL.}
\label{tab:estimand-separation}
\resizebox{\textwidth}{!}{%
\begin{tabular}{llrrrrrrr}
\toprule
Design & Method & $q_\mu^\star$ & $q_\eta^\star$ & $\widehat q$ & $F_{1,\mu}$ & $F_{1,\eta}$ & Exact support & ARI \\
\midrule
common $\kappa$ & M-CGL & 16 & 16 & 16.04 & 0.999 & 0.999 & 0.96 & 0.867 \\
common $\kappa$ & E-CGL & 16 & 16 & 16.06 & 0.998 & 0.998 & 0.95 & 0.867 \\
pure concentration & M-CGL & 0 & 16 & 0.96 & -- & 0.005 & 0.20 & 0.675 \\
pure concentration & E-CGL & 0 & 16 & 16.27 & -- & 0.992 & 0.83 & 0.633 \\
shared-natural-parameter & M-CGL & 100 & 20 & 21.13 & 0.348 & 0.974 & 0.00 & 0.859 \\
shared-natural-parameter & E-CGL & 100 & 20 & 20.02 & 0.333 & 1.000 & 0.98 & 0.870 \\
crossed support & M-CGL & 12 & 12 & 11.68 & 0.984 & 0.677 & 0.64 & 0.996 \\
crossed support & E-CGL & 12 & 12 & 11.89 & 0.681 & 0.980 & 0.69 & 0.996 \\
\bottomrule
\end{tabular}}
\end{table}

For the empty directional target, support $F_1$ is undefined and ``Exact
support'' is the probability of selecting the empty support. M-CGL selected
that support in 20\% of the pure-concentration repetitions, whereas E-CGL
recovered the nonempty
posterior-score contrast support with mean $F_{1,\eta}=0.992$. In the crossed design,
each method had high target-specific $F_1$ and lower $F_1$ against the other
target. The shared-natural-parameter design exposed a boundary for the directional
companion: its exact support was not selected in any replication. These results
motivate target-specific, rather than winner-take-all, comparisons.

\subsection{Sample-size recovery and oracle benchmarks}

Figure~\ref{fig:sample-size-trajectories} traces E-CGL as the sample size increases
from 300 to 2000. Centered-$\eta$ MSE decreased along every trajectory. At
$e_B=0.10$ with heterogeneous concentrations, the most difficult
small-sample sequence improved from $F_{1,\eta}=0.816$ and exact recovery 0.02
at $n=300$ to 0.999 and 0.97 at $n=2000$; the other trajectories approached
exact recovery at least as quickly.

\begin{figure}[t]
\centering
\includegraphics[width=\textwidth]{figure_03.pdf}
\caption{Sample-size trajectories for E-CGL. Points show Monte Carlo means,
and error bars show Monte Carlo standard errors over 100 replications. Both
oracle-error levels are shown under equal and heterogeneous concentrations.}
\label{fig:sample-size-trajectories}
\end{figure}

The path oracle is the fitted candidate with the largest $F_{1,\eta}$
within each repetition. The oracle-$S_\eta^\star$ estimator refits under the known
population support. We report
\[
\Delta_{\mathrm{sel}}
=F_{1,\eta}^{\mathrm{path\ oracle}}-F_{1,\eta}^{\mathrm{BIC}},
\qquad
\Delta_{\mathrm{NLL}}
=\mathrm{NLL}(\widehat\Theta)-
\mathrm{NLL}(\widehat\Theta_{S_\eta^\star}^{\mathrm{oracle\ support}}).
\]
These are finite-sample benchmarks, not asymptotic oracle claims. Selector
gaps were at most 0.0083 at $e_B=0.05$. The difficult heterogeneous-$\kappa$,
$e_B=0.10$ sequence had $\Delta_{\mathrm{sel}}=0.0478$ and
$\Delta_{\mathrm{NLL}}=0.1653$ at $n=300$, with the selector gap falling to
0.0009 by $n=2000$. At $n=300$, the separate path-oracle gap was 0.1357:
the exact support occurred on 6\% of fitted paths, and BIC selected it in
33.3\% of those covered repetitions. Thus candidate-path coverage, rather
than BIC selection alone, was the main small-sample limitation in this cell.

\subsection{Stress conditions and computation}

The five stress cells cover moderately and strongly dense support, $d>n$, and
weak beta-min. E-CGL attained $F_{1,\eta}=0.998$ in the sparse weak-signal cell
but 0.803 for $d=500,n=300$. With 80 active coordinates, $F_{1,\eta}$ improved
from 0.805 at $n=300$ to 0.901 at $n=1000$. E-ACGL reduced false positives in
some cells, with mixed effects on estimation and clustering; all results
appear in Table~\ref{tab:stress-results}.

\input{../simulations/tables/csda_main_section45_stress_260804.tex}

In a paired runtime diagnostic for the representative heterogeneous
$e_B=0.05$, $n=1000$, $d=200$ cell, median end-to-end wall-clock times and
interquartile ranges (IQRs) over 100 repetitions were 509.61
(415.62--597.18) seconds for M-CGL
and 74.62 (69.14--79.85) seconds for E-CGL. The sparse-method totals include the
paired dense initialization and the method-specific path, refit, and selector;
one-time Rcpp compilation was excluded. Because the methods use different path
and optimization rules, these values describe computational burden rather than
speed superiority.

\clearpage
\section{Real data analysis}

\subsection{Data and analysis protocol}

Classic3 contains abstracts from the CISI, CRAN, and MED document collections
and has been used in directional document clustering
\citep{banerjee2005vmf}. The raw payload contained
$n_0=3{,}890$ documents. We removed two exact duplicates and five prespecified
near-duplicates, leaving $n=3{,}883$ documents. The cleaned class counts were
1,454, 1,397, and 1,032, respectively.

Documents were encoded with the vocabulary-aligned Sparse Lexical and
Expansion Model (SPLADE), with a maximum input sequence length of 512
\citep{formal2021splade,formal2022hardnegative}. The exact encoding identifiers
were
\begingroup\small
\begin{quote}\noindent
Model: \texttt{naver/splade-cocondenser-ensembledistil}\\
Revision: \texttt{49cf4c7b0db5b870a401ddf5e2669993ef3699c7}.
\end{quote}
\endgroup
The payload and vocabulary checksums are reported in Supplementary Section
S9. Coordinates occurring in fewer than two documents were removed, the
remaining coordinates were ranked by variance over the unlabeled cleaned
payload, the top 2,000 were retained, and each vector was normalized to unit
Euclidean norm. The resulting matrix had no zero rows, duplicate
transformed rows, duplicate vocabulary entries, or nonfinite values. The
model identifier, revision, vocabulary hash, duplicate audit, and payload hash
were fixed before model fitting.

The preprocessing rules and split indices were prespecified. E-CGL uses the
same single 240-point path and explicit empty-support candidate as in the
primary simulations.

The benchmark resolution was fixed at $K=3$ to match the three supplied broad
topic labels. The labels were otherwise withheld from coordinate ranking,
initialization, parameter estimation, tuning, and support selection; only after
fitting were they used to calculate ARI and NMI and to name the aligned
components.
Five prespecified 80/20 splits provide stability and held-out diagnostics; they
are not pooled with the full-data estimates. After
the prespecified duplicate exclusions, each split ranks coordinates using its
training documents only. The resulting training-selected coordinate mask is
then applied without change to the test documents, after which both matrices
are row-normalized.

\subsection{Methods and evaluation}

The comparison comprised spherical $k$-means, permutation-tuned sparse
$k$-means \citep{witten2010sparseclustering}, dense vMF mixtures with shared
and component-specific concentrations, the Rossi-style M-L prototype selector,
M-CGL, E-CGL, and the secondary E-ACGL extension. Likelihood-based fits used
30 dense random starts.
The M-L path allowed at most 600 path points, including the zero-penalty fit.
E-CGL and E-ACGL used
the single zero-augmented 240-point relative path defined in Section~3; M-CGL
used the same 240-point dense-score-proxy path as in the numerical study.
Each sparse vMF method was followed by target-specific distinct-support
refitting and within-method BIC selection. M-CGL and the E-series methods also
included an explicit empty-support candidate; an empty prototype support is
incompatible with the unit-norm M-L directions. E-ACGL used the fixed weights
in \eqref{eq:eacgl-penalty}.
The same E-series rule was used in each prespecified Classic3 split.
The sparse vMF methods were selected by their method-specific BIC after
support-constrained refitting. Sparse $k$-means used a label-blind permutation
gap criterion with 25 permutations and 10 weight-bound candidates.

Selected $q$ is target-specific: weighted coordinates for sparse $k$-means,
prototype-union coordinates for M-L, directional-contrast coordinates for
M-CGL, and posterior-score contrast coordinates for E-CGL and E-ACGL.
Accordingly, differences in $q$ do not by themselves rank the estimands.
All methods were evaluated by ARI and NMI. Observed log-likelihood and the
practical nominal-dimension BIC were additionally reported for vMF models.

\subsection{Full-data results and interpretation}

The dense free-$\kappa_k$ mixture attained ARI 0.9760 and NMI 0.9538. E-CGL
selected 1,420 of 2,000 coordinates, excluding 29.00\% of the coordinates from
posterior-score contrasts while attaining ARI 0.9761 and NMI 0.9539.
M-CGL produced a closely matching partition (ARI 0.976 and NMI 0.954) with
1,471 directional-contrast coordinates.
M-L retained 1,942 prototype-support coordinates. E-ACGL selected 1,438
posterior-score contrast coordinates and attained ARI 0.9746 and NMI 0.9515.
The permutation-tuned sparse $k$-means fit selected 985 coordinates but had
ARI 0.1374; its selected weight bound was not at the evaluated grid boundary.
Independent restarts at that bound reproduced ARI 0.1373 and nearly the same
objective value, so this result is reported as performance of the
permutation-selected solution under the SPLADE representation rather than as
a single-start anomaly or a general statement about sparse $k$-means.
Table~\ref{tab:classic3-results} summarizes the full-data comparison.

\input{../simulations/tables/csda_table4_classic3_floor1e4_260812.tex}

For coordinate $j$, define
\[
R_j=\left(n^{-1}\sum_{i=1}^n x_{ij}^2\right)^{1/2},
\qquad
I_{kj}^{(\mathrm{ctr})}=R_j\left(\widehat\eta_{kj}-\widehat{\bar\eta}_j\right),
\qquad
I_j^{(\mathrm{com})}=R_j\left|\widehat{\bar\eta}_j\right|.
\]
After fitting, the E-CGL components were aligned to CISI, CRAN, and MED; their
concentrations were 738.32, 815.85, and 610.42. Panel A of
Fig.~\ref{fig:classic3-interpretation} selects the five largest positive
root-mean-square (RMS)-scaled contributions for each matched component and displays
their signed values $I_{kj}^{(\mathrm{ctr})}$. CISI contrasts were led
by \emph{library}, \emph{information}, and \emph{librarian}; CRAN by
\emph{flow}, \emph{mach}, and \emph{pressure}; and MED by \emph{tumor},
\emph{inhibitor}, and \emph{rat}. Panel B ranks inactive coordinates by
$I_j^{(\mathrm{com})}$ and shows those with the largest common
RMS-scaled natural-parameter contributions. Their fitted values are equal across
components, so their centered contrasts are numerically zero
even though the common natural-parameter baseline remains in the fitted
density.

\begin{figure}[!ht]
\centering
\includegraphics[width=\textwidth]{figure_04.pdf}
\caption{RMS-scaled posterior-score contrasts and common natural-parameter contributions in
Classic3. E-CGL selected 1,420 of 2,000 coordinates. Panel A shows signed
RMS-scaled posterior-score contrast contributions $I_{kj}^{(\mathrm{ctr})}$
(blue, positive; orange, negative). Panel B shows the ten inactive coordinates
with the largest RMS-scaled common contributions $I_j^{(\mathrm{com})}$. Their equal component values
represent retained common effects rather than posterior-score heterogeneity.
Supplied labels were used only after fitting to align and name the components.}
\label{fig:classic3-interpretation}
\end{figure}

\subsection{Classic3 split stability}

For each pair of the five prespecified Classic3 splits, the conditional support
Jaccard index is the support-intersection size divided by the support-union size
after both supports are restricted to the vocabulary shared by that pair.
E-CGL attained a mean of 0.933 over the ten split pairs, and the 15
selected-contrast tokens displayed in Fig.~\ref{fig:classic3-interpretation}
were retained in all five splits.
Because the true coordinate support is unknown, support TPR, FPR, precision,
and $F_1$ are not reported.

\section{Conclusions}

We introduced posterior-score contrast support for vMF mixtures and proposed
E-CGL to estimate it by penalizing centered natural-parameter contrasts while
preserving common baselines. This estimand identifies the coordinates entering
the feature-dependent linear terms of pairwise posterior log-odds. It differs
from prototype support and directional-contrast support, and the latter can
diverge from posterior-score contrast support under concentration
heterogeneity. The
guarded proximal GEM path, target-preserving refit, and within-method BIC
provide a practical computational procedure aligned with this target.

Across the evaluated $n=1000$ primary-grid cells, E-CGL selected 16.02--16.41
coordinates on average and attained $F_{1,\eta}$ between 0.988 and 0.999 for a
true support size of 16. At $n=1000$, support recovery changed little across
the evaluated oracle Bayes errors, although clustering agreement declined,
showing that clustering and
support difficulty need not move together. Recovery deteriorated when the same
contrast signal was distributed over many weak coordinates; in the moderately
dense $n=300$ condition, $F_{1,\eta}=0.805$ and ARI was 0.556.

In Classic3, E-CGL excluded 29.00\% of coordinates while attaining ARI and NMI
close to those of the dense component-specific-concentration fit. The five
prespecified splits gave a mean
conditional support Jaccard index of 0.933, and all 15 displayed contrast
tokens were retained in every split. These findings describe stability under
the fixed SPLADE representation and vocabulary filters, not recovery of an
unknown population support or invariance to arbitrary rotations of the feature
space.

The conclusions are conditional on fixed $K$ and on the path-generated model
collection. Refit BIC is used within method under the qualifications in
Section~\ref{subsec:bic-refit}; no general selection-consistency or global
optimality claim is made. The theoretical results concern a fixed-dimensional
conditional convex subproblem and an idealized residual-vanishing sequence, so
the high-dimensional experiments provide computational rather than asymptotic
evidence. Weak separation, diffuse signals, small samples, and local mixture
optima can still cause underselection, while adaptive weighting remains
sensitive to pilot contrasts and path coverage.

Further work includes effective degrees of freedom, joint uncertainty in $K$
and support, and selection under diffuse weak signals. E-CGL is therefore
intended as a target-specific sparse mixture procedure rather than a uniformly
superior replacement for dense or prototype-oriented models.

\appendix
\section{Proofs of theoretical results}
\label{app:proofs}

The main text states the analytic bounds that determine each result; this
appendix gives the complete arguments.

\begin{proof}[Proof of Lemma~\ref{lem:finite-existence}]
The closed simplex and the product of the $K$ closed Euclidean balls of radius
$\kappa_{\max}$ are compact. Each vMF component density is positive and
continuous on this set, so the observed log-likelihood is finite and
continuous. The centered group penalty is also continuous. The criterion
therefore attains its maximum by the extreme-value theorem.
\end{proof}

\begin{proof}[Proof of Proposition~\ref{prop:centered-support}]
Taking the average of $v=\bar v\mathbf1_K+c$ and using
$\mathbf1_K^\top c=0$ determines $\bar v$ and hence $c$ uniquely. The three
zero-contrast statements are immediate. Finally,
$\sum_{k<\ell}(v_k-v_\ell)^2=K\sum_k(v_k-\bar v)^2$, which gives the stated
identity.
\end{proof}

\begin{proof}[Proof of Proposition~\ref{prop:canonical-directional}]
Because $\kappa=\bar\kappa\mathbf1_K+h$ and
$M=\mathbf1_K\bar\mu^\top+U$,
\[
E=\operatorname{diag}(\kappa)M
=\bar\kappa\mathbf1_K\bar\mu^\top+\bar\kappa U
+h\bar\mu^\top+\operatorname{diag}(h)U.
\]
Apply $H_K$, using $H_K\mathbf1_K=0$, $H_KU=U$, and $H_Kh=h$.
\end{proof}

\begin{proof}[Proof of Proposition~\ref{prop:label-invariance}]
Because $P\mathbf1_K=\mathbf1_K$, $P$ commutes with $J_K$ and $H_K$.
Permutation matrices are orthogonal and therefore preserve the Euclidean
norm. Applying Proposition~\ref{prop:centered-prox} gives equivariance of the
proximal map.
\end{proof}

\begin{proof}[Proof of Proposition~\ref{prop:centered-prox}]
The spaces spanned by $\mathbf1_K$ and orthogonal to $\mathbf1_K$ are
orthogonal. The penalty is zero on the first space and equals the Euclidean
norm on the second. The minimization therefore preserves $J_Kv$ and applies
the standard group soft-thresholding map to $H_Kv$.
\end{proof}

\begin{proof}[Proof of Proposition~\ref{prop:majorization}]
The vMF log-partition function has Hessian equal to the covariance matrix of a
unit-norm random vector and is therefore positive semidefinite with operator
norm at most one. The Hessian of the $k$th block of $f_t$ consequently has
operator norm at most $N_k$. The descent result follows by combining the
smooth majorization inequality with optimality of the proximal update.
\end{proof}

\begin{proof}[Proof of Theorem~\ref{thm:fixed-responsibility}]
For finite natural parameters, the vMF log-partition function is strictly
convex. Its Hessian is the covariance matrix of a vMF random vector, and the
vMF density is positive on the full sphere; hence
$\operatorname{Var}(a^\top X)>0$ for every nonzero $a\in\mathbb R^d$.
Since $N_k>0$, the sum of its $K$ blocks is strictly convex. The standard
large-argument expansion of the modified Bessel function gives

\[
-\log C_d(\kappa)
=\kappa-\frac{d-1}{2}\log\kappa+R_d(\kappa),
\qquad \kappa\longrightarrow\infty.
\]

Since $R_d(\kappa)=O(1)$, there exist $\kappa_0$ and $B_d<\infty$ such that
$R_d(\kappa)\geq-B_d$ for every $\kappa\geq\kappa_0$.

Consequently, as $\lVert\eta_k\rVert_2\to\infty$,

\[
\begin{aligned}
N_k\{-\log C_d(\lVert\eta_k\rVert_2)\}
-r_k^\top\eta_k
\geq{}&
(N_k-\lVert r_k\rVert_2)\lVert\eta_k\rVert_2
\\
&-\frac{N_k(d-1)}{2}\log\lVert\eta_k\rVert_2
-N_kB_d.
\end{aligned}
\]

The positive linear coefficient implied by
$\lVert r_k\rVert_2<N_k$ dominates the logarithmic term. Thus $f_t$ is
coercive. Adding the finite convex group penalty preserves strict convexity and
coercivity, so $F_t^{(a)}$ has a unique minimizer. If that minimizer is interior to
the concentration bound, it is also the unique bounded-subproblem solution.
The gradient of $f_t$ is globally Lipschitz with constant no larger than
$\max_kN_k$. Hence the initial trial step $1/\max_kN_k$ satisfies the exact
majorization inequality and is bounded away from zero. Standard
proximal-gradient theory yields convergence and the $O(u^{-1})$ objective rate
\citep{beck2009fista}. The strict inequality
$\lVert r_k\rVert_2<N_k$ excludes the degenerate case in which all
observations receiving positive responsibility from a component have one
common direction. The finite cap still gives existence in that case, but
this theorem does not apply.
\end{proof}

\begin{proof}[Proof of Proposition~\ref{prop:gem-monotonicity}]
The standard EM lower-bound identity expresses the observed log-likelihood
difference as the auxiliary-function difference plus a nonnegative
Kullback--Leibler divergence term. The penalty depends only on the parameters
and can be subtracted from both sides, yielding the penalized result.
\end{proof}

\begin{proof}[Proof of Theorem~\ref{thm:idealized-gem}]
Let $\Theta^{(t_j)}\to\bar\Theta$ be any convergent subsequence, which exists
by compactness. Asymptotic regularity also gives
$\Theta^{(t_j+1)}\to\bar\Theta$. After passage to a further subsequence if
needed, boundedness of the step sizes gives $s_{t_j}\to\bar s$ for some
$\bar s\in[\underline s,\overline s]$. Exact posterior responsibilities are
continuous on the assumed compact interior set, so
$\tau^{(t_j)}\to\tau(\bar\Theta)$.

The gradient $\nabla f_\tau(E)$ is jointly continuous in $(E,\tau)$ on this
set. For $s$ bounded away from zero, the scaled proximal map is jointly
continuous in its input and in $s$. Moreover,
$s_{t_j}\mathcal G_{s_{t_j}}
(E^{(t_j+1)};\tau^{(t_j)})\to0$. The vanishing natural-parameter residual
therefore implies

\[
\bar E
=\operatorname{prox}_{\bar s\lambda_\eta\mathcal P_a}
\{\bar E-\bar s\nabla f_{\tau(\bar\Theta)}(\bar E)\}.
\]

The proximal fixed-point characterization yields

\[
0
\in
\nabla f_{\tau(\bar\Theta)}(\bar E)
+\lambda_\eta\partial\mathcal P_a(\bar E).
\]

At exact posterior responsibilities, Fisher's identity for the observed-data
score gives
$\nabla f_{\tau(\bar\Theta)}(\bar E)=-\nabla_E\ell(\bar\Theta)$, proving the
natural-parameter KKT inclusion. For the mixing block,

\[
\pi^{(t_j+1)}-\frac{N(\tau^{(t_j)})}{n}\longrightarrow0,
\]

while $\pi^{(t_j+1)}\to\bar\pi$ and
$N(\tau^{(t_j)})\to N\{\tau(\bar\Theta)\}$. Hence
$\bar\pi_k=N_k\{\tau(\bar\Theta)\}/n$. Because the limiting mixing proportions are
strictly positive,

\[
\frac{\partial\ell(\bar\Theta)}{\partial\pi_k}
=\frac{N_k\{\tau(\bar\Theta)\}}{\bar\pi_k}
=n
\]

for every $k$, which is the interior simplex KKT condition with $\nu=n$.
The objective, residuals, and KKT system are equivariant under any fixed
component permutation. The argument does not apply to boundary or
finite-tolerance cases excluded by the assumptions
\citep{green1990penalizedem,wu1983em,tseng2004emproximal}.
\end{proof}

\section*{Data and code availability}

The source corpora must be obtained from their original providers under the
applicable terms. Derived document-level SPLADE matrices are not redistributed
because redistribution permission for the source corpora has not been
established. A redistribution-safe reproducibility archive containing the
reviewed implementation, tests, preprocessing specifications, configuration
files, random seeds and split records, exclusion records, vocabulary and
payload checksums, compact numerical summaries, and result-generation scripts
is available from the corresponding author on reasonable request. Large raw,
path, cache, and private project files are not part of that archive.

\section*{Funding}

This research received no specific grant from funding agencies in the public,
commercial, or not-for-profit sectors.

\section*{Declaration of generative AI and AI-assisted technologies in the manuscript preparation process}

During the preparation of this work, the authors used OpenAI ChatGPT and Codex
for English-language editing, translation, software-development assistance,
and consistency checking. The authors reviewed and verified the resulting
text, code, references, and numerical claims and take full responsibility for
the content of the publication.

\bibliographystyle{elsarticle-harv}
\bibliography{references}
\end{document}
